\documentclass[10pt,a4paper]{amsart}
\usepackage[margin=19mm]{geometry}
\usepackage{amsmath,amssymb,mathtools}
\usepackage[scr=rsfs,cal=euler]{mathalfa}
\usepackage{xfrac}
\usepackage[unicode,hidelinks,bookmarks=false]{hyperref}
\newcommand{\ie}{i.e.\ }
\newcommand{\cf}{cf.\ }
\newcommand{\resp}{respectively\ }
\newcommand{\Subsection}{\subsection}
\providecommand{\nobreakdash}{\nobreak\hbox{-}}
\setlength{\emergencystretch}{2em}
\multlinegap0pt
\usepackage{amssymb}
\newtheorem{theorem}{Theorem}
\def \T{\mathbb{T}}
\newcommand{\cle}{\mathcal{E}}
\newcommand{\cls}{\mathcal{S}}
\newcommand{\vp}{\varphi}
\title{Norm attaining dual truncated Toeplitz operators}
\author{Sudip Ranjan Bhuia, Puspendu Nag}
\date{Source: arXiv:2601.09375v2 (CC BY 4.0)}
\begin{document}
\maketitle

\noindent\emph{Source and licence.} Norm attaining dual truncated Toeplitz operators. Version arXiv:2601.09375v2 (CC BY 4.0), distributed by arXiv under the Creative Commons Attribution 4.0 International licence, http://creativecommons.org/licenses/by/4.0/, as recorded in the arXiv licence metadata for this exact version and shown on the abstract page https://arxiv.org/abs/2601.09375v2. Full licence notice: \texttt{licenses/arxiv-2601.09375-CC-BY-4.0.txt}.\par\smallskip

\noindent\textit{Editorial scope.} Counted unit: the article's theorem thm:NA-density-finite-Blaschke (source0.tex lines 2371--2382). The article's complete original proof of it is delivered with it (source0.tex lines 2384--2474, one \texttt{proof} environment of 91 lines). No mathematics is written, repaired or re-derived: the statement and the proof are the article's own lines, in the article's own notation, and no retained line is substituted or re-flowed.  Retained uncounted as the source-local context the proof needs: lines 2205--2218.  Nothing was omitted from any retained span, so the package carries no omission record.  No retained line contains a citation, so the wrapper carries no bibliography and no bibliography entry.  No retained line contains a figure environment, an image-inclusion command or drawing code, so no image is delivered and the package declares no asset dependency.  Editorial formatting only, in addition: an AMS \texttt{amsart} preamble that restates the article's own declarations and macros used by the retained lines, namely the environments \texttt{theorem} on separate counters, together with the standard packages those lines need.  The retained environments print in this document as Theorem 1 in order of appearance, whereas the article prints the same items with the numbering of its own full document.  The complete native source of the article as distributed in the arXiv e-print for this exact version is bundled unchanged as \texttt{sources/192668/source0.tex}.
\begin{theorem}\label{thm:finite-Blaschke-norm-attainment}
Let $u$ be a finite Blaschke product and let $\vp\in L^\infty$. Then
\[
D_\vp\text{ is norm attaining}\quad\Longleftrightarrow\quad
m\left(\{|\vp|=\|\vp\|_\infty\}\right)>0.
\]
If $\vp\neq0$ and $u$ has degree $n$, then whenever $D_\vp$ is norm
attaining,
\[
\operatorname{codim}_{L^2(E_\vp)}\cle_\vp=\dim\cls_{\vp,u}
\leq 2n.
\]
In particular, $\cle_\vp$ is infinite-dimensional.
\end{theorem}

\begin{theorem}[Density of norm attaining and non-norm attaining DTTOs]\label{thm:NA-density-finite-Blaschke}
Let $u$ be a finite Blaschke product. Then both sets
\[
\mathcal{N}_u:=\{\vp\in L^\infty:D_\vp\in\mathcal{NA}\}
\qquad\text{and}\qquad L^\infty\setminus\mathcal{N}_u
=\{\vp\in L^\infty:D_\vp\notin\mathcal{NA}\}
\]
are norm dense in $L^\infty$.
	
Consequently, both the norm attaining and the non-norm attaining
dual truncated Toeplitz operators are operator-norm dense in the class $\{D_\vp:\vp\in L^\infty\}$.
\end{theorem}

\begin{proof}
By Theorem~\ref{thm:finite-Blaschke-norm-attainment},
\[
D_\vp\in\mathcal{NA}\quad\Longleftrightarrow\quad m(E_\vp)>0,
\qquad E_\vp:=\{\zeta\in\T:|\vp(\zeta)|=\|\vp\|_\infty\}.
\]
	
We first prove that $\mathcal{N}_u$ is dense in $L^\infty$. Let $\vp\in L^\infty$ and $\varepsilon>0$. If $D_\vp\in\mathcal{NA}$, we may simply take $\psi=\vp$. Suppose, therefore, that $D_\vp\notin\mathcal{NA}$, and set $r:=\|\vp\|_\infty$. Since $D_0=0$ is norm attaining, necessarily $r>0$.
	
Set
\[
\eta:=\min\left\{\frac{\varepsilon}{2},\frac{r}{2}\right\}.
\]
By the definition of the essential supremum, the set $A_\eta:=\{\zeta\in\T:|\vp(\zeta)|>r-\eta\}$ has positive measure.
Since $\eta<r$, we also have $\vp\neq0$ a.e.\ on $A_\eta$. Define
\[
\psi(\zeta)
=\begin{cases}
		r\dfrac{\vp(\zeta)}{|\vp(\zeta)|},
		& \zeta\in A_\eta,\\[3mm]
		\vp(\zeta),
		& \zeta\notin A_\eta.
\end{cases}
\]
Then $\psi\in L^\infty$ and $\|\psi\|_\infty=r$. Moreover, for $\zeta\in A_\eta$,
\[
|\psi(\zeta)-\vp(\zeta)|=r-|\vp(\zeta)|<\eta,
\]
whereas $\psi=\vp$ on $\T\setminus A_\eta$. Hence $\|\psi-\vp\|_\infty\leq\eta<\varepsilon$.
	
Furthermore, $|\psi|=r$ a.e.\ on $A_\eta$, so $m(E_\psi)\geq m(A_\eta)>0$. Therefore, Theorem~\ref{thm:finite-Blaschke-norm-attainment} implies that
$D_\psi$ is norm attaining. Thus $\mathcal{N}_u$ is norm dense in
$L^\infty$.
	
We now prove that $L^\infty\setminus\mathcal{N}_u$ is norm dense in
$L^\infty$. Let $\vp\in L^\infty$ and $\varepsilon>0$. If $D_\vp\notin\mathcal{NA}$, we again take $\psi=\vp$. Suppose that
$D_\vp\in\mathcal{NA}$.
	
First assume that $\vp\neq0$, and set $r:=\|\vp\|_\infty>0$. By Theorem~\ref{thm:finite-Blaschke-norm-attainment}, $m(E_\vp)>0$. Since Lebesgue measure on $\T$ is nonatomic, there exists a measurable function $h:E_\vp\to(0,1]$ such that $\operatorname*{ess\,inf}_{E_\vp}h=0$. For instance, one may choose
pairwise disjoint measurable sets $E_n\subseteq E_\vp$ of positive
measure, define $h=1/n$ on $E_n$, and define $h=1$ on the remaining
part of $E_\vp$.
	
Choose
\[
0<\delta<\min\left\{1,\frac{\varepsilon}{r}\right\},
\]
and define
\[
\psi(\zeta)=
\begin{cases}
(1-\delta h(\zeta))\vp(\zeta),
& \zeta\in E_\vp,\\
\vp(\zeta),
& \zeta\notin E_\vp.
\end{cases}
\]
Then $\|\psi-\vp\|_\infty\leq\delta r<\varepsilon$.
	
We claim that $\|\psi\|_\infty=r$. Clearly $|\psi|\leq r$ a.e.
On the other hand, since $\operatorname*{ess\,inf}_{E_\vp}h=0$, for every $\eta>0$ the set $\{\zeta\in E_\vp:h(\zeta)<\eta\}$ has positive measure. On this set,
\[
|\psi(\zeta)|=r(1-\delta h(\zeta))>r(1-\delta\eta).
\]
Hence $\|\psi\|_\infty\geq r(1-\delta\eta)$ for every $\eta>0$.
Letting $\eta\downarrow0$ yields $\|\psi\|_\infty\geq r$, and therefore $\|\psi\|_\infty=r$.
	
However, on $E_\vp$ we have $h>0$ a.e., so $|\psi|=r(1-\delta h)<r$ a.e. On $\T\setminus E_\vp$, $|\psi|=|\vp|<r$ a.e. Consequently,
$|\psi|<r=\|\psi\|_\infty$ a.e.\ on $\T$, and hence $m(E_\psi)=0$. By
Theorem~\ref{thm:finite-Blaschke-norm-attainment}, $D_\psi$ is not norm attaining.
	
It remains to consider $\vp=0$. Given $\varepsilon>0$, choose
$0<c<\varepsilon$ and define
\[
\psi(z):=c\frac{1+z}{2},\, z\in \T.
\]
Then $\|\psi\|_\infty=c<\varepsilon$, while, for $z=e^{it}$,
\[
|\psi(e^{it})|=c\left|\cos\frac{t}{2}\right|<c
\]
for almost every $t$. Thus $m(E_\psi)=0$, and Theorem~\ref{thm:finite-Blaschke-norm-attainment} implies that
$D_\psi$ is not norm attaining. Therefore, $L^\infty\setminus\mathcal{N}_u$ is norm dense in $L^\infty$.
	
Finally, for $\vp,\psi\in L^\infty$, we have
\[
\|D_\vp-D_\psi\|=\|D_{\vp-\psi}\|=\|\vp-\psi\|_\infty.
\]
Thus the map $\vp\mapsto D_\vp$ is isometric. Hence the two density
statements above translate directly into operator-norm density of the
norm attaining and non-norm attaining DTTOs in the class $\{D_\vp:\vp\in L^\infty\}$.
\end{proof}

\end{document}
